\section{Related Work}
\label{sec:related}

This section discusses state-of-the-art approaches categorized by filtering, anomaly detection,
provenance tracking, and kernel-hardening, highlighting their strengths and limitations compared
to the DeSFAM framework. Table~\ref{tab:related_work} summarizes recent syscall filtering and
anomaly detection approaches, highlighting their problem focus, limitations, and core techniques.

\begin{table*}[tp]
  \centering
  \caption{Comparison of State-of-the-Art Container Security Solutions.}
  \label{tab:related_work}
  \includegraphics[width=\linewidth]{table1}
\end{table*}

\subsection{System Call Filtering Approaches}
\label{subsec:syscall_filtering}

Confine~\cite{ref31_confine} uses static analysis to construct restrictive seccomp profiles,
achieving a 47.6\% syscall reduction rate (SRR). However, its static nature prevents adaptation
to runtime behaviors, reducing effectiveness against evolving workloads. Similarly,
Prof-gen \cite{ref36_profgen} relies on static techniques that either require manual
configuration or omit dynamic tracing, limiting scalability and robustness in real-world
deployments.

KubeRosy~\cite{ref29_kuberosy} and Optimus~\cite{ref27_optimus} incorporate eBPF for live
syscall monitoring and runtime policy updates. KubeRosy leverages Kubernetes CRDs for real-time
policy refinement but remains tied to Kubernetes environments. Optimus employs
association-based filtering to iteratively adapt syscall profiles, yet its automation is
constrained---it requires operator approval for unmatched syscalls and halts profile updates if
association analysis fails. Prof-gen, while combining dynamic boot-time tracing with static call
graph analysis to generate minimal access lists without source code, lacks support for capturing
behavioral shifts during runtime execution.

\subsection{Anomaly-Based Intrusion Detection Systems}
\label{subsec:anomaly_ids}

Machine learning-based techniques have been widely explored for syscall anomaly detection. The
BoSC-based IDS employs Bag-of-System Calls modeling with classifiers like Random Forests,
offering reasonable accuracy in small-scale environments. However, frequent retraining
requirements hinder scalability and adaptability.

Lightweight models such as the TinyML-based Escape Detector target IoT and edge devices but fall
short in handling complex attack vectors in high-performance cloud systems due to limited
adaptability and accuracy.

DeSFAM addresses these limitations in Phase~2 via a hybrid anomaly detector combining a
Variational Autoencoder (VAE) with an Isolation Forest, trained on the DongTing
dataset~\cite{ref26_dongtingdataset}. This approach enables accurate, low-latency detection of
zero-day attacks with no need for retraining, while supporting real-time sliding window
inference.

\subsection{Provenance-Based Detection Approaches}
\label{subsec:provenance}

Provenance-driven systems like PACED monitor syscall relationships and causality across
containers to detect cross-namespace attacks. While effective, they suffer from high overhead
due to continuous graph construction and limited scalability.

DeSFAM circumvents this by embedding syscall context (e.g., PID, namespace, timestamps)
directly into model features, enabling behavioral awareness without the cost of full provenance
tracking.

\subsection{Kernel-Level Security Enhancements}
\label{subsec:kernel_security}

Kernel-layer defenses such as Condo use Trusted Execution Environments (TEEs) to safeguard
permission integrity. Though secure, their reliance on specialized hardware and kernel
modification makes them unsuitable for large-scale cloud deployments.

In contrast, DeSFAM achieves portability and scalability by utilizing eBPF maps and LSM hooks to
enforce syscall policies dynamically, without modifying the kernel or depending on TEEs.

\subsection{Common Limitations of Existing Approaches}
\label{subsec:common_limitations}

While there have been meaningful advances in syscall minimization and container security, current
solutions face several key limitations:

\begin{itemize}
  \item Static policies cannot adapt to evolving runtime behaviors.
  \item Lack of real-time enforcement allows syscall abuse between updates.
  \item Machine learning approaches require frequent retraining, limiting scalability.
  \item Hardware dependencies (e.g., TEEs) restrict deployment flexibility.
\end{itemize}

\subsection{DeSFAM: Advancing the State of the Art}
\label{subsec:desfam_advancement}

DeSFAM addresses key limitations of prior systems through a unified, multi-phase architecture.
It combines hybrid syscall profiling, combining static binary disassembly with dynamic runtime
tracing, to generate accurate syscall access lists. Real-time enforcement is achieved via eBPF
probes and LSM hooks, enabling low-latency blocking and policy updates without restarting
containers. Anomaly detection is handled by SyscallAD, a hybrid model using a Variational
Autoencoder and Isolation Forest for precise, retraining-free detection of known and zero-day
attacks. DeSFAM also incorporates CVE-aware filtering and MITRE ATT\&CK-based risk scoring to
restrict vulnerable syscalls dynamically, and applies tiered response logic to enforce
fine-grained actions based on threat severity. The framework operates without kernel modification
or reliance on orchestration platforms, offering scalable and adaptive syscall security for
modern container environments.
